% Замените подсказки в квадратных скобках содержанием своей задачи.
% Сборка: latexmk -pdf -outdir=output 02-миниобзор-шаблон.tex
% Перед сборкой убедитесь, что output/ указан в .gitignore.
\documentclass[12pt,a4paper]{article}
\usepackage[T2A]{fontenc}
\usepackage[utf8]{inputenc}
\usepackage[russian]{babel}
\usepackage{amsmath,amssymb,amsthm}
\usepackage[a4paper,margin=2.5cm]{geometry}
\usepackage[hidelinks]{hyperref}
\newtheorem{lemma}{Лемма}

\title{Миниобзор: [тема]}
\author{[имя]}
\date{\today}

\begin{document}
\maketitle

\section{Объекты и примеры}
[Определите объекты, которые встречаются в выбранной теме. Разберите
простейший пример и укажите, какие вопросы он вызывает.]

\section{Основные публикации и результаты}
[Свяжите две-три ключевые работы в короткий рассказ. Для каждого результата
укажите точные условия, заключение и номер теоремы. Проверьте это по самой
публикации. Формат ссылки:]

\noindent\verb|\cite[теорема 2.3]{source1}|

\section{Граница известного и гипотезы}
[Разделите доказанный результат, вопрос авторов и собственную гипотезу.
Если вопрос назван открытым в статье, укажите версию статьи и дату проверки
позднейших работ. Запишите один пример или граничный случай.]

\section{Дальнейшие вопросы}
[Запишите вопросы, возникшие после чтения и примеров. Укажите, что ещё
следует проверить по литературе или на новых примерах.]

\begin{thebibliography}{9}
\bibitem{source1}
[Авторы], \emph{[Название]}, [издание или arXiv], [год], [постоянная ссылка].
\end{thebibliography}

\end{document}
